\vspace{-0.1em}
La formation stellaire dans les galaxies est alimentée par le gaz du milieu interstellaire (MIS), qui comprend des phases moléculaire, atomique et ionisée. Cette étude examine la relation entre la formation stellaire et les différentes phases du MIS dans un échantillon de 16 galaxies proches, dont 9 naines/de type tardif et 7 spirales.
Nous avons utilisé les données CO de \textit{HERACLES}, hydrog\`ene atomique (\HI) de \textit{THINGS}, à 24~$\mu$m de \textit{SINGS} et UV de \textit{GALEX NGS}. Nous avons décomposé les spectres de l'\HI\ à l'aide de l'algorithme ROHSA, qui prend en compte la cohérence spatiale des spectres voisins, afin de séparer les composantes \HI\ cinématiquement étroites et larges, que nous utilisons comme traceurs du milieu neutre froid (CNM) et  (WNM), plutôt que comme des mesures directes des phases thermiques.\\
Nous examinons la répartition spatiale des composantes cinématiquement froide et tiède et constatons qu'elles s'étendent toutes deux jusqu'aux régions les plus externes des disques atomiques. La composante étroite représente en médiane environ 45\% de la masse totale d'\HI\ et paraît plus filamentaire que la composante large. Sa fraction n'est pas systématiquement plus élevée dans les bras spiraux, contrairement aux attentes théoriques selon lesquelles le refroidissement dans les bras favoriserait la formation de gaz froid. De plus, nous ne trouvons pas d'indication claire que la composante étroite soit généralement plus concentrée près des émissions CO que la composante large, comme on pourrait s'y attendre si elle était plus étroitement associée au gaz moléculaire.\\
Nous examinons la relation de Kennicutt--Schmidt pour les différentes composantes gazeuses afin de déterminer lesquelles sont les plus étroitement liées à la formation stellaire. Dans les naines, la composante cinématiquement froide présente une corrélation globale plus forte avec le taux de formation stellaire que l'\HI\ total ou la composante tiède. Dans les spirales, c'est la composante tiède qui est globalement mieux corrélée. Cette différence pourrait s'expliquer par le fait que la largeur des raies renseigne sur des conditions gazeuses différentes dans ces deux types de galaxies. Dans les naines, l'émission étroite trace probablement du gaz plus froid et moins turbulent, favorable à la formation stellaire. Dans les spirales, les régions de formation stellaire peuvent contenir du gaz froid dont les raies sont élargies par la turbulence ou des mouvements non résolus. Globalement, le gaz moléculaire reste la composante la plus étroitement liée au taux de formation stellaire, comme l'ont montré les études antérieures. Le rôle du gaz atomique cinématiquement froid dans la transition \HI--H$_2$ reste ambigu à notre résolution sub-kiloparsec.\\
\noindent\textbf{Mots-clés :} Galaxies proches, Milieu interstellaire, Formation stellaire
\newpage